\scanpage{165}
\mspage{477}
(\textsuperscript{108}2) (6 octobre) Dans ce ``déployement de force'' il n'y a nulle intention ``agressive'' au sens courant du terme, consciente ni inconsciente, seulement un désir inconscient d'impressionner, de forcer l'estime. Il est vrai que ce terme ``forcer l'estime'' qui me revient spontanément, porte déjà une connotation de c\,o\,n\,t\,r\,a\,i\,n\,t\,e, voisine de celle ``d'agression''. Cette intention inconsciente de contrainte, perçue au niveau inconscient également, doit souvent être vécu\unsure{e} comme une sorte d'agression (alors même que ce vécu reste occulté, tout comme le\unsure{s} réactions antagonistes qu'il déclenche). En même temps, il doit se faire souvent l'amalgame de ce vécu avec des vécus analogues, remontant à l'enfance, avec le père comme protagoniste, et où celui-ci apparaît comme le principal détenteur d'autorité répressive, voire comme un rival écrasant, envié et détesté.

Même sans un tel amalgame, et indépendamment aussi de toute perception che\unsure{z} autrui d'une intention de ``contrainte'' en moi, il doit y avoir souvent la perception d'un fort d\,é\,s\,é\,q\,u\,i\,l\,i\,b\,r\,e, d'une disharmonie foncière, dans ce ``déployement de force'' exclusivement yang (dans son esprit et intention, au moins). Cette démesure est néfaste au principal intéressé, savoir moi-même, et à la limite bel et bien ``dangereux'' pour sa survie physique même (comme des incidents-santé de ces dernières années me l'ont bien montré! ). C'est cela sans doute qui a été en filigrane dans ma pensée, quand j'ecrivais qu' ``un tel déployeme\unsure{nt} de force'' était ressenti ``en tout cas comme dangereux'' - dangereux ``par nature'', un exemple donc à surtout ne pas suivre...! Un tel ressenti est sûrement suffisant pour susciter des ``réactions de défense'', en l'absence même de toute agression ou intention d'agresser.

Il est vrai que de telles relations d'ambiguité se sont reproduites après 1976, avec certains de mes élèves notamment, en des moments où tout investissement mathématique était absent, et où il n'y avait aucun ``dépoyement de force'' apparent dans ma vie. Il est vrai aussi que les `` déployements'' en question du p\,a\,s\,s\,é ont créé une r\,é\,p\,u\,t\,a\,t\,i\,o\,n, qui continue à me coller à la peau, surtout dans ma vie professionelle, et qui dans une certaine mesure se substitue à la perception de celui que je suis d\,a\,n\,s\quad l\,e\quad p\,r\,é\,s\,e\,n\,t. De plus, j'ai acquis dans le commerce de certains thèmes mathématiques une aisance telle que, même en dehors de mes périodes mathématiques et ma réputation aidant, cette aisance ou maîtrise naturelle peuvent avoir déjà l'effet du ``déployement de force'', sur des élèves peu motivés, et me faire ressentir par eux (en dépit de certains traits avenants voire rassurants) comme une sorte de Superman (un peu Superpère sur les bords!).\margin{\unsure{-}}
% archive sheet number: 164
